\scanpage{106}

La forme plus faible, ``géométrique'' de la conjecture ci-dessus dirait que \deleted{\unsure{c'est}} l'on obtient bien ainsi, par le schéma associé à $M^\circ\otimes_ZQ$ (sur lequel $\pi$ opère ...), la \deleted{\unsure{groupe de}} ``partie compacte'' du \deleted{\unsure{son}} groupe de Galois motivique de $\overline{\mathbb F}_p$, qui serait donc isom. à
\[
G\simeq\uline{D}(\unsure{M_0|M^\circ}\otimes Q)\times G_m
\]
avec l'hom inj: $G_m\longrightarrow G$ s'identifiant à $j$, et \deleted{$\unsure{G_m}$} $(x,\lambda)\longrightarrow\lambda^2$ s'identifie à $\varepsilon$.

En termes de valeurs propres de Frobenius:

\uline{Conjecture faible} \unsure{Soit|Les} $\lambda\in\bar Q^*$, \deleted{\unsure{tel qu'il}\inserted{\deleted{\unsure{tel}}} \unsure{existe un entier $n>0$ avec $\lambda^n$ satisfaisant aux conditions, de la forme $p$}} un \uline{entier} algébrique, \deleted{tel qu'il existe un entier} tel qu'il existe un entier $i\geq0$, de sorte que $\lambda=p^{i/2}u$, \deleted{\unsure{soit un}} où $u$ est une \uline{unité absolue} \uline{exc.} $p$, et qu'il existe un entier $n>0$ tel que $u^n\in Q'(i)^*$, et $v_{\mathfrak p}(u^n)+v_{\mathfrak p'}(u^n)=0$ pour tt couple d'idéaux premiers de $Q'(i)$ induisant le même idéal premier de $Q'$. Alors il existe un entier $m>0$ tel que $\lambda^m=(p^m)^{i/2}u^m$ soit \inserted{une} valeur propre de frobenius \unsure{$f_{q^m}$} dans un motif semi-simple \inserted{de poids $i$} convenable sur $\mathbb F_{q^m}$

\margin{Résulte \unsure{probablement} de la théorie des multiplications complexes: cf. \uline{Taniyama--Shimura}}

\mspage{105}
